% !PRINTED 246
\weberpage{246}

Wenn eine der beiden Grössen $a_1$, $a_2$ verschwindet, so haben wir, da dann nicht alle Wurzeln von gleichem Zeichen sein können, eine oder zwei positive Wurzeln.

Wenn $a_3 = 0$ ist, so reducirt sich die Discriminante auf
\[
a_2^2(a_1^2 - 4a_2),
\]
und wenn diese positiv ist, so hat die cubische Gleichung eine verschwindende und noch zwei andere reelle Wurzeln. Diese sind positiv, wenn
\[
a_1 < 0, \quad a_2 > 0,
\]
negativ, wenn
\[
a_1 > 0, \quad a_2 > 0,
\]
und es ist eine von ihnen positiv und eine negativ, wenn
\[
a_1 \gtreqless 0, \quad a_2 < 0
\]
ist.

\begin{center}
\S.\,78.

\gesperrt{Discussion der biquadratischen Gleichung.}
\end{center}

Bei den Gleichungen vierten und fünften Grades existiren entweder keine oder zwei oder vier imaginäre Wurzeln. Ist die Discriminante negativ, so hat man zwei imaginäre Wurzeln. Bei positiver Discriminante können entweder vier oder keine imaginären Wurzeln vorhanden sein. Diese beiden Fälle zu unterscheiden, wird weiterhin unsere Aufgabe sein. Wir wenden uns aber zunächst zu einer elementaren Betrachtung der biquadratischen Gleichung, die wir der Einfachheit halber in der verkürzten Form
\weq{1}{x^4 + ax^2 + bx + c = 0}
annehmen wollen, von der wir leicht zur allgemeinen Form zurückkehren können.

Bezeichnen wir die Wurzeln mit $\alpha$, $\beta$, $\gamma$, $\delta$ und setzen, wie in \S.\,47
\weq{2}{\begin{aligned} \alpha - \beta &= v+w, & \gamma - \delta &= v-w,\\ \alpha - \gamma &= w+u, & \delta - \beta &= w-u,\\ \alpha - \delta &= u+v, & \beta - \gamma &= u-v, \end{aligned}}
so sind $u^2$, $v^2$, $w^2$ die Wurzeln der cubischen Resolvente
\weq{3}{y^3 + 2ay^2 + (a^2-4c)y - b^2 = 0,}
und die Discriminante $D$ dieser cubischen Gleichung, die durch
%CONT>>
